\documentclass{butexFR}


\titre{Solution aux besoins}
\UE{Analyse}
\sujet{SAÉ 3.02 - Labyrinthe}
\enseignant{Mme Boneva}
\eleves{Bredelle Victor \\ Marouze Antonin \\ Harlaut Romain \\ Zheng Angèl \\ Lavogiez Baptiste}


\begin{document}

\fairepagedegarde
\fairemarges
\tabledematieres



\section{Introduction}

Ce rapport présente l'analyse UML du projet \textbf{SAÉ S3.02 - Labyrinthe}, un jeu avec deux modes (libre et progression) :

\begin{enumerate}
    \item Diagrammes de cas d'utilisation
    \item Fiches descriptives
    \item Prototype d'interface
    \item Diagramme de classe
\end{enumerate}

\section{Équipe et organisation du travail}
\textbf{Groupe G3.3}

\begin{itemize}
    \item Bredelle Victor
    \item Marouze Antonin
    \item Harlaut Romain
    \item Zheng Angèl
    \item Lavogiez Baptiste
\end{itemize}

\tsec{Victor}

Rédaction de trois fiches descriptives avec Baptiste : Charger profil, Afficher progression, Lancer mode progression.

\tsec{Antonin}

Réalisation de la maquette \textit{Figma} avec Angèl et Romain.

\tsec{Romain}

Réalisation de la fiche descriptive "Déplacer joueur" et aide à la maquette \textit{Figma}.

\tsec{Angèl}

Réalisation de la maquette \textit{Figma} avec Antonin et Romain.

\tsec{Baptiste}

Rédaction de trois fiches descriptives avec Victor et rédaction du rapport.

\tsecnonum{Le groupe}

\info{\textbf{L'ensemble des contributeurs a collaboré à l'élaboration du diagramme de classes.}}

\section{Cas d'utilisation}

\subsection{Diagramme de cas d'utilisation}

\begin{figure}[h!]
    \centering
    \includegraphics[width=0.5\textwidth]{figures/cas_dutilisation/diagramme_de_cas_dutilisation/default/fig_1.png}
    \caption{Cas d'utilisation identifiés}
    \label{fig:cas_dutilisation_diagramme_de_cas_dutilisation_1}
\end{figure}

\subsection{Analyse du diagramme de cas d'utilisation}

Le diagramme identifie \textbf{sept cas d'utilisation} répartis entre deux acteurs.

\subsubsection{Acteurs identifiés}

\begin{itemize}
    \item \textbf{Utilisateur} : Acteur principal représentant le joueur.
    \item \textbf{Système} : Acteur secondaire représentant les processus de persistance automatique des données.
\end{itemize}

L'acteur "Système" modélise la sauvegarde automatique post-action.

\subsubsection{Description des cas d'utilisation}

Le diagramme couvre le \textbf{mode progression} :

\begin{enumerate}
    \item \textbf{Jouer} : Point d'entrée du jeu.
    \item \textbf{Lancer en mode progression} : Initialisation du mode progression (saisie du nom).
    \item \textbf{Charger profil} : Reprise d'une progression existante.
    \item \textbf{Commencer une partie} : Démarrage d'un labyrinthe (défi ou configuration libre).
    \item \textbf{Afficher la progression} : Consultation du tableau de bord du joueur.
    \item \textbf{Se déplacer dans le labyrinthe} : Déplacement case par case.
    \item \textbf{Sauvegarder la progression} : Persistance automatique de l'état du joueur.
\end{enumerate}

\subsubsection{Clarifications}

Choix de modélisation :

\begin{itemize}
    \item \textbf{"Charger profil"} spécialise "Lancer mode progression" si le nom existe déjà.

    \item \textbf{Sauvegarde automatique} : Modélisée par l'acteur "Système" pour distinguer l'automatisme des actions explicites.

    \item \textbf{Mode libre} : Réutilise "Jouer" et "Commencer une partie" (avec configuration manuelle : largeur, hauteur, pourcentage de murs) mais sans authentification, profil ni sauvegarde. Le joueur accède directement au labyrinthe après configuration.
\end{itemize}

\newpage

\subsubsection{Cas d'utilisation : Lancer mode progression}
\textbf{Système:} Labyrinthe\\
\textbf{Acteur principal:} Utilisateur\\
\textbf{Déclencheur:} Aucun\\
\textbf{Autres acteurs:} Aucun\\
\textbf{Préconditions:} Aucun\\
\textbf{Garanties en cas de succès:} L'utilisateur accède au mode progression.\\
\textbf{Garanties minimales:} L'intégrité de la base de données est préservée.\\
\textbf{Scénario nominal:}
\begin{enumerate}
    \item L'utilisateur sélectionne ``Lancer mode progression''.
    \item Le système affiche un menu où l'utilisateur peut saisir son nom.
    \item L'utilisateur saisit son nom.
    \item Le système vérifie que le nom est valide, sauvegarde son profil dans la base de données et lance le mode progression.
\end{enumerate}
\textbf{Scénarios alternatifs:}

\emph{A: Le système détecte une erreur lors de la vérification}\\
\hspace*{1em}4(A): Le système affiche une notification d'erreur, puis revient à l'étape 2.

\emph{B: Le système détecte une occurrence du nom dans la base de données (charger profil)}\\
\hspace*{1em}4(B): Le système démarre le \textbf{Cas d'utilisation} ``Charger profil''.

\newpage

\subsection{Fiches descriptives}
\subsubsection{Cas d'utilisation : Charger profil}
\textbf{Système:} Labyrinthe\\
\textbf{Acteur principal:} Utilisateur\\
\textbf{Déclencheur:} Le système détecte une occurence du nom dans la base de données.\\
\textbf{Préconditions:} L'utilisateur est dans le menu de saisie du nom.\\
\textbf{Garanties en cas de succès:} Le système charge le profil associé au nom et l'utilisateur reprend sa progression.\\
\textbf{Garanties minimales:} L'intégrité des données existantes est préservée.\\
\textbf{Scénario nominal:}
\begin{enumerate}
    \item L'utilisateur a rentré un nom déjà existant lors du cas ``Lancer mode progression''.
    \item Le système affiche un menu où l'utilisateur peut choisir entre "Nouveau profil" ou "Reprendre".
    \item L'utilisateur sélectionne "Reprendre".
    \item Le système charge le profil associé au nom et l'utilisateur reprend sa progression.
\end{enumerate}
\textbf{Scénarios alternatifs:}

\emph{A: L'utilisateur choisit "Nouveau profil"}\\
\hspace*{1em}2(A): Le système écrase les données du profil existant et crée un nouveau profil vide. L'utilisateur accède alors au mode progression.

\tsecnonum{Clarification}

Ce cas est la suite du scénario alternatif B de "Lancer mode progression".

\newpage

\subsubsection{Cas d'utilisation : Afficher progression}
\textbf{Système:} Labyrinthe\\
\textbf{Acteur principal:} Utilisateur\\
\textbf{Déclencheur:} Aucun\\
\textbf{Autres acteurs:} Aucun\\
\textbf{Préconditions:} L'utilisateur possède un profil et se trouve en mode progression.\\
\textbf{Garanties en cas de succès:} Le système affiche la progression associée au joueur.\\
\textbf{Garanties minimales:} Les données affichées sont cohérentes et non modifiées.\\
\textbf{Scénario nominal:}
\begin{enumerate}
    \item L'utilisateur sélectionne la fonctionnalité ``accéder au menu du mode progression''.
    \item Le système calcule le score et les caractéristiques du joueur, puis affiche sa progression.
\end{enumerate}
\textbf{Scénarios alternatifs:} Aucun

\newpage

\subsubsection{Cas d'utilisation : Déplacer joueur}
\textbf{Système:} Labyrinthe\\
\textbf{Acteur principal:} Utilisateur\\
\textbf{Déclencheur:} Aucun\\
\textbf{Autres acteurs:} Aucun\\
\textbf{Préconditions:} Le joueur se trouve dans le labyrinthe.\\
\textbf{Garanties en cas de succès:} Le système déplace le joueur à la case souhaitée.\\
\textbf{Garanties minimales:} La position du joueur reste cohérente.\\
\textbf{Scénario nominal:}
\begin{enumerate}
    \item Le joueur appuie sur une flèche directionnelle.
    \item Le système vérifie que la direction est accessible et déplace le joueur sur la case associée.
\end{enumerate}
\textbf{Scénarios alternatifs:}

\emph{A: Le joueur atteint la sortie}\\
\hspace*{1em}2(A): Le système notifie la réussite, sauvegarde la progression et redirige vers le menu de sélection de défi.

\emph{B: Le joueur se déplace dans une case inaccessible}\\
\hspace*{1em}2(B): Le système ne modifie pas la position du joueur et le scénario reprend à l'étape 1.

\newpage

\subsection{Prototypes d'interface}


\begin{figure}[h!]
    \centering
    \includegraphics[width=0.9\textwidth]{figures/cas_dutilisation/prototypes_dinterface/default/fig_1.png}
    \caption{Prototypes d'interface identifiés avec arborescence}
    \label{fig:cas_dutilisation_prototypes_dinterface_1}
\end{figure}

\smallskip

\subsubsection{Vue d'ensemble et parcours utilisateur}

Le prototype couvre \textbf{huit écrans} répartis entre les deux modes de jeu.

\textbf{Parcours utilisateur -- Mode Progression (6 écrans) :}

\begin{enumerate}
    \item \textbf{Menu d'accueil} : Options "Jouer" et "Quitter".
    \item \textbf{Choix du mode} : Sélection "Mode Libre" ou "Mode Progression".
    \item \textbf{Saisie du nom} : Identification ou création du profil.
    \item \textbf{Gestion du profil existant} : Choix "Nouvelle partie" ou "Reprendre".
    \item \textbf{Sélection étape/défi} : Tableau de progression avec code couleur.
    \item \textbf{Interface de jeu} : Labyrinthe avec joueur et sortie.
\end{enumerate}

\textbf{Parcours utilisateur -- Mode Libre (4 écrans) :}

\begin{enumerate}
    \item \textbf{Menu d'accueil}
    \item \textbf{Choix du mode}
    \item \textbf{Configuration du labyrinthe} : Paramétrage des dimensions et murs.
    \item \textbf{Interface de jeu} : Labyrinthe généré aléatoirement.
\end{enumerate}

\subsubsection{Correspondance avec les cas d'utilisation}

Le tableau suivant établit la correspondance entre les écrans du prototype (mode progression) et les cas d'utilisation.

\vspace{1em}
\begin{center}
\begin{tabular}{|p{5cm}|p{8cm}|}
\hline
\textbf{Cas d'utilisation} & \textbf{Écrans concernés} \\
\hline
Lancer mode progression & Choix du mode $\rightarrow$ Saisie du nom \\
\hline
Charger profil & Page de profil existant (choix "Reprendre" ou "Nouvelle partie") \\
\hline
Jouer & Menu d'accueil $\rightarrow$ Clic sur jouer\\
\hline
Commencer une partie & Sélection étape/défi (clic sur un défi) $\rightarrow$ Labyrinthe en jeu \\
\hline
Afficher progression & Sélection étape/défi (affichage et survol des défis) \\
\hline
Se déplacer dans le labyrinthe & Labyrinthe en jeu \\
\hline
Sauvegarder progression & Automatique après sortie du labyrinthe \\
\hline
\end{tabular}
\end{center}
\vspace{1em}

Cette traçabilité assure la correspondance entre fonctionnalités et interface.

\subsubsection{Clarifications}

Boutons "Retour"/"Quitter" omis par clarté.

L'écran de progression s'adapte au profil chargé.

Cet écran intègre visualisation et sélection des défis. Le survol affiche les détails, le clic lance le défi.

Détail :

\smallskip


\begin{figure}[h!]
    \centering
    \includegraphics[width=1\textwidth]{figures/cas_dutilisation/prototypes_dinterface/clarifications/fig_2.png}
    \caption{Page de progression avec survol d'un défi}
    \label{fig:cas_dutilisation_prototypes_dinterface_2}
\end{figure}

À la sortie du labyrinthe, sauvegarde et mise à jour automatiques.

\newpage

\section{Diagramme de classe}


\begin{figure}[h!]
    \centering
    \includegraphics[width=0.9\textwidth]{figures/diagramme_de_classe/default/default/fig_1.png}
    \caption{Diagramme de classe identifié}
    \label{fig:diagramme_de_classe_default_1}
\end{figure}


\subsection{Structure générale}

L'architecture se décompose en trois niveaux conceptuels :

\begin{enumerate}
\item \textbf{Niveau contrôleur} : \texttt{Main} orchestre le flux général.
\item \textbf{Niveau logique métier} : \texttt{GameMode}, \texttt{Maze}, \texttt{Player}, etc. encapsulent les règles métier.
\item \textbf{Niveau persistance} : \texttt{PlayerDatabase} isole les opérations de sauvegarde/chargement.
\end{enumerate}

Cette séparation facilite la maintenance et la testabilité.

\subsection{Conformité avec la méthode vue en cours}

La méthodologie préconise une classe \texttt{Main} avec une fonction par cas d'utilisation, structure respectée ici.

\subsection{Correspondance avec les cas d'utilisation}

La correspondance suivante démontre que la classe \texttt{Main} couvre l'ensemble des cas d'utilisation.

\vspace{1em}
\begin{center}
\begin{tabular}{|p{5cm}|p{8cm}|}
\hline
\textbf{Fonctions du \texttt{Main}} & \textbf{Cas d'utilisation correspondants} \\
\hline
\texttt{play()} & Jouer \\
\hline
\texttt{startGame(): GameMode} & Lancer mode progression\\
\hline
\texttt{movePlayerPosition(...)} & Se déplacer dans le labyrinthe \\
\hline
\texttt{displayProgress(player: Player)} & Afficher progression \\
\hline
\texttt{savePlayer(player: Player)} & Sauvegarder progression\\
\hline
\texttt{loadPlayer(name: String): Player} & Charger profil \\
\hline
\texttt{createMaze(...)} & Commencer une partie\\
\hline
\end{tabular}
\end{center}
\vspace{1em}

Le type abstrait \texttt{GameMode} mutualise le code commun aux deux modes.

\subsection{Clarifications}

\tsecnonum{Les étapes et défis}

\info{Un \texttt{Level} correspond à une \textit{Étape} (pour éviter le conflit avec \texttt{Stage} de JavaFX). Un \texttt{Challenge} est un défi.}

Il y a 3 étapes et 3 défis par étape, les structures de données ont donc une taille fixe.

Un objet \texttt{Challenge} comporte les attributs requis par le sujet.

La méthode \texttt{isCompleted(): boolean} de \texttt{Level} retourne \texttt{True} si au moins un défi est complété.

Les méthodes d'accès direct aux étapes (getters/setters) ne sont pas mentionnées car non demandées dans le plan.

\tsecnonum{Le type \texttt{GameMode}}

\texttt{startGame()} retourne un \texttt{GameMode} (\texttt{FreeMode} ou \texttt{ProgressionMode}). Ce type abstrait unifie les opérations communes et évite la duplication.

\begin{figure}[h!]
\centering
\includegraphics[width=0.8\textwidth]{figures/diagramme_de_classe/default/default/fig_2.png}
\caption{startGame() dans le Menu}
\label{fig:diagramme_de_classe_default_2}
\end{figure}

\tsecnonum{L'initialisation de \texttt{Player} dans \texttt{ProgressionMode}}
Le constructeur de \texttt{ProgressionMode} invoque \texttt{loadPlayer(name: String): Player} pour :
\begin{itemize}
\item Charger un profil existant si le nom est trouvé.
\item Créer un nouveau profil sinon.
\end{itemize}

Le \texttt{Player} devient attribut de \texttt{ProgressionMode}, garantissant l'exclusivité de gestion des profils.

\section{Conclusion}

Cette analyse respecte les besoins du jalon 1 et établit une base solide pour l'implémentation des deux premières étapes.

\end{document}

